\section{Introduction}\label{sec:introduction}

Prescribing multiple points of a plane curve is an interpolation problem:
how small can its degree be when the positions and the required orders of
vanishing are given? Nagata's conjecture predicts the answer at the
square-root scale for sufficiently many very general points. Nagata
introduced the problem in his work on Hilbert's fourteenth problem and
proved the square cases $r\ge16$~\cite{Nagata1959,Nagata1965}. We prove
the strict inequality for every $r\ge10$, allowing unequal multiplicities.

We first work over $\C$. Let
\[
 U_r=\{(p_1,\ldots,p_r)\in(\Pp^2_\C)^r:p_i\ne p_j\text{ for }i\ne j\}.
\]
A nonzero effective plane curve is the zero divisor of a nonzero
homogeneous form of degree $d\ge1$. Its multiplicity at a point is the
order of that form in the regular local ring. Components may repeat.
The phrase \emph{very general} in the following theorem means the
complement of one countable union of proper closed subsets, chosen
simultaneously for all degrees and multiplicity vectors.

The scale can already be seen when all prescribed multiplicities equal
$m$. A degree-$d$ form has $\binom{d+2}{2}$ coefficients, and order at
least $m$ at $r$ points imposes at most $r\binom{m+1}{2}$ linear
conditions. Their leading terms balance at $d=m\sqrt r$. A count of
conditions supplies curves asymptotically above this scale; proving their
absence below it requires control of the dependencies among the conditions, for every
$m$. The strict endpoint also matters: at nine points a nonzero cubic
always exists, so the asserted strict inequality cannot start at $r=9$.

\begin{theorem}\label{thm:main}
For every integer $r\ge10$, there is a countable union $E_r$ of proper
Zariski-closed subsets of $U_r$, with nonempty complement, such that the
following holds for every tuple $(p_1,\ldots,p_r)\notin E_r$.
If $C$ is a nonzero effective plane curve of degree $d\ge1$ and
$m_1,\ldots,m_r\in\Z_{\ge0}$ satisfy $\mult_{p_i}C\ge m_i$, then
\[
 d\sqrt r>\sum_{i=1}^r m_i.
\]
\end{theorem}

On the blowup of the plane at such a tuple, the theorem gives a real divisor
class of self-intersection zero that intersects every curve positively.
It also determines the multipoint Seshadri constant: the largest $t$ for
which the pullback of a line minus $t$ times the sum of the exceptional
curves has nonnegative intersection with every curve. Its value is
$1/\sqrt r$, which is irrational when $r$ is nonsquare.
Section~\ref{sec:consequences} proves these consequences
and extends the theorem to every uncountable algebraically closed field
of characteristic zero.

\subsection{Prior work and the interpolation approach}

The geometry of plane linear systems connects this problem with the cones
of effective and nef divisor classes on blowups. Ciliberto--Harbourne--
Miranda--Ro\'e~\cite{CHMR2012} develop these connections and compare
several forms and extensions of Nagata's conjecture.
For homogeneous multiplicities, Ro\'e~\cite{Roe2001} obtained the
uniform estimate $d>m(\sqrt{r-1}-\pi/8)$ for $r\ge10$ by
specializing to configurations with infinitely near points.
Harbourne~\cite{Harbourne2001} sharpened rational lower bounds and
proved further cases in restricted multiplicity ranges.
For ten very general points, degeneration methods give
$d/m\ge117/37$ for every nonempty homogeneous system: Eckl's Seshadri estimate and
Ciliberto--Dumitrescu--Miranda--Ro\'e's analysis of plane systems
give this bound in their respective formulations~\cite{Eckl2011,CDMR2011}.
Such bounds explain the importance of controlling arbitrarily large
multiplicities at each fixed nonsquare $r$.

Dumnicki--Harbourne--K\"uronya--Ro\'e--Szemberg~\cite{DHKRS2016}
study an analogue for very general monomial valuations, while
Nivoche~\cite{Nivoche2019} relates the conjecture to transcendental
interpolation in $\C^n$. These perspectives emphasize that the obstruction
is not merely the number of imposed conditions, but how vanishing
conditions interact with the geometry of the ambient surface.
More recently, Ciliberto--Miranda--Ro\'e~\cite[Theorem~1]{CMR2026}
constructed irrational square-zero nef rays on the blowups at every
$r\ge10$ very general points. Those constructions do not determine the
equal-weight multipoint Seshadri constant. Nagata's conjecture singles out
the class with equal exceptional coefficients, whose nefness gives that
constant's exact value.

The proof below uses two established approaches to this interaction.
Ciliberto--Miranda~\cite{CM1998} develop degenerations of plane linear
systems, and Evain~\cite{Evain2007} studies their infinitesimal limits
when fat points, meaning points equipped with specified orders of
vanishing, approach a divisor. Dumnicki~\cite{Dumnicki2007} encodes
interpolation by finite diagrams of monomial exponents and their
evaluation matrices. In particular, the proof of his Proposition~13 converts a
one-point derivative matrix into polynomial evaluation on the exponent
diagram. Our final rank problem has the same evaluation interpretation,
with a different set of allowed polynomial degrees supplied by a
collinear collision.

The elliptic stage uses the classical multiplicative theta product and
its residue-class Laurent bases; the conventions used here are recorded
in Rosengren~\cite[arXiv version, Section~1.2 and Exercise~1.6.5]{Rosengren2020}.
We prove the needed product identities, basis formulas and convergence
estimates in the normalization required by the argument.

\subsection{Proof strategy}

Closed incidence loci turn a failure at very general points into a
system of curves existing at every configuration. A product over cyclic
permutations makes the multiplicities equal; write $d$ and $m$ for the
resulting degree and common multiplicity. Moving the points toward a
smooth curve gives a nonzero polynomial on its normal line bundle,
with order at least $m$ at arbitrary normal displacements. For square
$r$, its two highest coefficients already exclude $d/m\le\sqrt r$,
including equality (Section~\ref{sec:reductions}).

For nonsquare $r$, we use a cubic and set nine normal displacements to
zero. Dividing out the forced coefficient zeros and changing the fiber
coordinate give a finite-dimensional space of polynomial sections whose
first $m+1$ coefficient bundles have equal degree
(Section~\ref{sec:elliptic}). On each fixed elliptic curve, the remaining
$q=r-9$ points collide along a horizontal coordinate line. After $b<m$
fiber derivatives, a section must then have base order at least $q(m-b)$;
thus a particular finite jet map has a nonzero kernel. A separate
degeneration gives explicit normalized theta bases whose expressions
in logarithmic base and fiber coordinates $(x,y)$ tend to $e^{Kx+jy}$.
Here $j$ is the fiber degree and $K$ a Laurent exponent. Their mixed
derivatives give the limiting matrix
\[
 \bigl(K^\ell j^b\bigr),\qquad
 0\le b<m,\quad 0\le\ell<q(m-b),
\]
with columns indexed by exponent pairs $(j,K)$
(Section~\ref{sec:limits}). Only these numerical matrices are compared
across curves; the argument does not choose a simultaneous family of
kernel vectors.

To exclude the kernel, a total count of exponent pairs is not enough.
We group them into horizontal rows of fixed $K$, with $j$ varying.
Interpolation one row at a time gives full column rank if no row has
more than $m$ points and at most $q(m-t+1)$ rows have at least $t$ points.
The exponent pairs lie in a polygon whose slice-width profile supplies
exactly these bounds. Since $d/m$ is rational and $r$ is nonsquare,
the strict gap $d/m<\sqrt r$ permits a contraction of that polygon;
taking a common power of the original curve absorbs the integer endpoint
error in the row count. Full rank persists for small positive degeneration
parameter, contradicting the kernel on each fixed curve
(Section~\ref{sec:interpolation}).
